\documentclass[12pt,a4paper]{article}

\usepackage[T1]{fontenc}
\usepackage[utf8]{inputenc}
\usepackage{microtype}
\usepackage[margin=2.5cm]{geometry}
\usepackage{hyperref}
\usepackage{xcolor}
\usepackage{abstract}
\usepackage{array}
\usepackage{booktabs}
\usepackage{graphicx}
\usepackage{amsmath}
\usepackage{amssymb}
\usepackage{float}
\usepackage{listings}
\usepackage{tabularx}

\hypersetup{
  colorlinks=true,
  linkcolor=blue!60!black,
  urlcolor=blue!60!black,
  citecolor=blue!60!black
}

\lstdefinestyle{aspirelisting}{
  basicstyle=\footnotesize\ttfamily,
  breaklines=true,
  breakatwhitespace=false,
  frame=single,
  rulecolor=\color{black!25},
  backgroundcolor=\color{black!2},
  framesep=3mm,
  xleftmargin=0.75em,
  xrightmargin=0.25em,
  columns=fullflexible,
  keepspaces=true,
  showstringspaces=false,
  aboveskip=0.8em,
  belowskip=0.8em
}

\lstset{style=aspirelisting}

\renewcommand{\abstractnamefont}{\normalfont\bfseries}
\renewcommand{\abstracttextfont}{\normalfont\small}

\title{Orchestrating PostgreSQL with .NET Aspire:\\ A Practical Guide to AppHost, pgAdmin,\\ and Persistent Containers in Existing .NET Solutions}
\author{K.C. Voss}
\date{2026}

\begin{document}
\maketitle

\begin{abstract}
\noindent
.NET Aspire 9 provides a composable orchestration layer that wires containerised infrastructure --- PostgreSQL, pgAdmin, Redis, and more --- into existing .NET solutions with a handful of builder calls and zero manual connection-string management. This article documents, in step-by-step and canonical-code form, how to add an Aspire AppHost project to any existing .NET~9 solution: declaring a PostgreSQL container with a persistent named Docker volume, attaching a pgAdmin~4 browser interface, enforcing \texttt{ContainerLifetime.Persistent} so containers survive successive \texttt{dotnet run} calls, wiring a logical database to the web project via \texttt{WithReference} and \texttt{WaitFor}, and publishing OpenTelemetry signals through a shared ServiceDefaults library. All generated secrets are stored in .NET User Secrets; no credentials reach source control. The worked example uses the HobbitGame solution --- a Blazor Server application backed by PostgreSQL --- to ground every instruction in a concrete, reproducible context.

\vspace{0.5em}
\noindent\textbf{Keywords:} .NET Aspire; PostgreSQL; pgAdmin; Docker; AppHost; ServiceDefaults; OpenTelemetry; Blazor; Entity Framework Core; container orchestration; developer experience.
\end{abstract}

\tableofcontents
\bigskip

\section{Introduction}
\label{sec:introduction}

The canonical path for connecting a .NET web application to a PostgreSQL database has historically involved a chain of manual steps: installing the database engine, creating a database and user, writing a connection string into \texttt{appsettings.json}, and remembering to exclude credentials from source control. Each step is individually simple; in aggregate they represent a setup tax that compounds across team members, CI pipelines, and deployment environments. When the application is containerised the tax increases: the developer must now also author Docker Compose files, manage container lifetimes, and ensure the application does not start before the database is ready.

.NET Aspire~9 eliminates this tax at the local development layer. An AppHost project --- a thin \texttt{.exe} that hosts a \texttt{DistributedApplication} --- becomes the single entry point for the entire development stack. It declares containers as typed resources, wires connection strings into dependent projects at startup, enforces readiness ordering via \texttt{WaitFor}, and exposes a browser-based dashboard with aggregated logs, metrics, and traces. Postgres passwords are generated randomly on first run and stored in .NET User Secrets; they are never written to \texttt{appsettings.json} or committed to source control.

This article documents every concrete step required to add that orchestration layer to an existing .NET~9 solution. Section~\ref{sec:background} reviews the architectural primitives that Aspire introduces. Section~\ref{sec:apphost} covers the creation and structure of the AppHost project. Section~\ref{sec:wiring} details the canonical \texttt{Program.cs} resource wiring. Section~\ref{sec:servicedefaults} introduces the ServiceDefaults shared library and its observability configuration. Section~\ref{sec:web} describes the changes required in the orchestrated web project. Section~\ref{sec:secrets} explains the User Secrets flow and common credential pitfalls. Section~\ref{sec:discussion} addresses limitations and production generalisation. Section~\ref{sec:conclusion} concludes.

Throughout, the HobbitGame solution --- a Blazor Server application where an LLM-driven hobbit wanders a map of the Shire and persists state in PostgreSQL --- serves as the worked example, providing concrete project names, database resource names, and code excerpts that can be adapted to any other solution.

\section{Background: .NET Aspire Primitives}
\label{sec:background}

\subsection{The DistributedApplication Model}

.NET Aspire introduces a \texttt{DistributedApplication} abstraction whose sole responsibility is declaring \textit{resources} --- projects, containers, databases, message buses --- and the \textit{relationships} between them. At runtime the AppHost process orchestrates the lifecycle of every declared resource: it pulls container images, creates Docker volumes, injects environment variables and connection strings into dependent processes, and surfaces a unified health-and-telemetry dashboard~\cite{aspiredocs2024}.

The model separates three concerns that are conflated in traditional \texttt{docker-compose.yml} files: (i) \textit{infrastructure topology} (what exists and what depends on what), (ii) \textit{configuration injection} (how connection strings reach application code), and (iii) \textit{observability routing} (where logs, metrics, and traces are collected). Aspire handles all three through typed C\# builder calls rather than YAML, making the setup amenable to refactoring, code review, and static analysis.

\subsection{Resource Types Relevant to PostgreSQL Setups}

Three resource types are relevant to the PostgreSQL + pgAdmin scenario described in this article:

\begin{itemize}
  \item \textbf{Container resources} --- wrapping a Docker image. \texttt{AddPostgres} and \texttt{WithPgAdmin} both produce container resources. Aspire manages their start order, port allocation, and environment variables automatically.
  \item \textbf{Database resources} --- logical databases registered inside a container resource via \texttt{AddDatabase}. A database resource carries the connection string that will be injected into dependent projects.
  \item \textbf{Project resources} --- references to other \texttt{.csproj} files in the solution, declared via \texttt{AddProject\textless T\textgreater}. Aspire compiles and launches these projects as child processes, injecting the connection strings of any referenced database resources as \texttt{ConnectionStrings\_\_\{name\}} environment variables.
\end{itemize}

\subsection{Container Lifetime Options}

By default, Aspire tears down container resources when the AppHost process exits. For development scenarios where database state must survive restarts --- such as a game save file stored in PostgreSQL --- this default is undesirable. \texttt{ContainerLifetime.Persistent} instructs Aspire to leave the container running after the AppHost exits and to reuse it on the next invocation~\cite{aspiredocs2024}. The container is identified by the resource name; a name collision with an existing container will reuse rather than recreate it. Named Docker volumes behave analogously: \texttt{WithDataVolume()} creates or reuses a named volume, so rows written in one session are available in the next.

\subsection{Readiness Ordering with WaitFor}

PostgreSQL containers require a few seconds to initialise. If the web application starts before the database is ready, connection attempts fail and --- in the absence of retry logic --- the application either crashes or enters a degraded state. Aspire's \texttt{WaitFor} call on the project resource instructs the AppHost to delay process launch until the referenced resource has passed its health check. For Postgres, the health check queries the server's readiness endpoint; the web app process does not start until that check passes~\cite{aspiredocs2024}.

\subsection{ServiceDefaults and OpenTelemetry}

The ServiceDefaults library is a class library project (tagged \texttt{IsAspireSharedProject}) that centralises three cross-cutting concerns: OpenTelemetry export (traces, metrics, and logs via OTLP to the Aspire dashboard), HTTP client resilience defaults (\texttt{Microsoft.Extensions.Http.Resilience}), and service discovery (\texttt{Microsoft.Extensions.ServiceDiscovery}). Every orchestrated project calls \texttt{builder.AddServiceDefaults()} early in its \texttt{Program.cs}, giving the Aspire dashboard visibility into the full distributed trace~\cite{aspiredocs2024,otelspec2023}.

\section{Creating the AppHost Project}
\label{sec:apphost}

\subsection{Project Structure and the Secondary SDK}

The AppHost project is a conventional \texttt{.exe} project with one unusual property: it carries a \textit{secondary SDK} declaration that injects the Aspire build tooling. The secondary SDK line must appear immediately after the opening \texttt{<Project>} tag and before the \texttt{<PropertyGroup>}:

\begin{lstlisting}[language=XML]
<!-- src/HobbitGame.AppHost/HobbitGame.AppHost.csproj -->
<Project Sdk="Microsoft.NET.Sdk">
  <Sdk Name="Aspire.AppHost.Sdk" Version="9.0.0" />

  <PropertyGroup>
    <OutputType>Exe</OutputType>
    <TargetFramework>net9.0</TargetFramework>
    <ImplicitUsings>enable</ImplicitUsings>
    <Nullable>enable</Nullable>
    <IsAspireHost>true</IsAspireHost>
    <UserSecretsId>hobbit-apphost-secrets</UserSecretsId>
  </PropertyGroup>

  <ItemGroup>
    <PackageReference Include="Aspire.Hosting.AppHost"    Version="9.0.0" />
    <PackageReference Include="Aspire.Hosting.PostgreSQL" Version="9.0.0" />
  </ItemGroup>

  <ItemGroup>
    <ProjectReference Include="..\HobbitGame.Web\HobbitGame.Web.csproj" />
  </ItemGroup>
</Project>
\end{lstlisting}

Four properties in this file merit explicit comment.

\textbf{\texttt{Aspire.AppHost.Sdk}.} This is not the primary SDK (\texttt{Microsoft.NET.Sdk}); it is a secondary SDK that augments the build with Aspire-specific MSBuild targets. Without it, the \texttt{IsAspireHost} flag and the \texttt{Projects.*} type generation described in Section~\ref{sec:wiring} will not function.

\textbf{\texttt{IsAspireHost}.} Marks the project as the Aspire orchestrator. This flag enables the code-generation step that produces the \texttt{Projects.HobbitGame\_Web} type reference used in \texttt{Program.cs}.

\textbf{\texttt{UserSecretsId}.} This identifier links the project to a .NET User Secrets store. Aspire generates PostgreSQL credentials on first run and writes them to this store. Without a \texttt{UserSecretsId}, credential storage fails silently and the password is lost between sessions~\cite{microsoftusersecrets2024}.

\textbf{\texttt{Aspire.Hosting.PostgreSQL}.} This NuGet package adds the \texttt{AddPostgres}, \texttt{WithPgAdmin}, and \texttt{AddDatabase} extension methods to the builder. Without it, those calls do not exist.

\subsection{Project Reference Requirement}

The AppHost must carry a \texttt{ProjectReference} to every project it orchestrates. This requirement exists at two levels: (i) the build system must be able to compile the referenced project so that its output assembly is available for the AppHost to launch; (ii) the Aspire SDK code-generation step must be able to emit the \texttt{Projects.*} type that is used in the \texttt{AddProject\textless T\textgreater} call. A missing reference produces a build error of the form \texttt{CS0246: The type or namespace name `Projects' could not be found}.

\subsection{Adding the Project to the Solution}

After creating the \texttt{.csproj} file, the project must be registered in the solution file. The conventional command is:

\begin{lstlisting}
dotnet sln HobbitGame.sln add src/HobbitGame.AppHost/HobbitGame.AppHost.csproj
\end{lstlisting}

The same command applies to the ServiceDefaults project described in Section~\ref{sec:servicedefaults}. Visual Studio and Rider will both correctly resolve the secondary SDK and recognise the project as an Aspire orchestrator after the solution is reloaded.

\subsection{Setting the Startup Project}

Once the AppHost project is in the solution, it should be set as the startup project. In Visual Studio this is done by right-clicking the project in Solution Explorer and selecting \textit{Set as Startup Project}. In Rider, the run configuration dropdown should point at \texttt{HobbitGame.AppHost}. Running any other project directly will start only that process in isolation, without Aspire orchestration, and the database connection string will not be injected.

\section{Resource Wiring in Program.cs}
\label{sec:wiring}

\subsection{Canonical Structure}

The complete canonical file for the Game scenario is:

\begin{lstlisting}[language={[Sharp]C}]
// src/Game.AppHost/Program.cs
var builder = DistributedApplication.CreateBuilder(args);

var postgres = builder.AddPostgres("postgres")
    .WithDataVolume(isReadOnly: false)
    .WithPgAdmin()
    .WithLifetime(ContainerLifetime.Persistent);

var gameDb = postgres.AddDatabase("gamedb");

builder.AddProject<Projects.Game_Web>("webapp")
    .WithReference(gameDb)
    .WaitFor(gameDb)
    .WithExternalHttpEndpoints();

builder.Build().Run();
\end{lstlisting}

Each builder call in this file is load-bearing. Table~\ref{tab:builder-calls} documents the responsibility of every call, the failure mode produced by its omission, and the Aspire component that implements it.

\begin{table}[H]
\centering
\caption{Builder calls in AppHost \texttt{Program.cs} and the consequences of omitting each.}
\label{tab:builder-calls}
\small
\renewcommand{\arraystretch}{1.15}
\begin{tabularx}{\textwidth}{@{}>{\raggedright\arraybackslash}p{4.5cm} >{\raggedright\arraybackslash}X >{\raggedright\arraybackslash}p{3.5cm}@{}}
\toprule
\textbf{Call} & \textbf{Effect} & \textbf{Omission failure mode} \\
\midrule
\texttt{AddPostgres("postgres")} & Pulls \texttt{postgres} Docker image; assigns resource name used as DNS alias & No Postgres container; web app cannot connect \\
\addlinespace
\texttt{.WithDataVolume()} & Creates/reuses named Docker volume & Database state lost on container stop \\
\addlinespace
\texttt{.WithPgAdmin()} & Pulls \texttt{dpage/pgadmin4}; auto-registers the server connection & No browser DB UI \\
\addlinespace
\texttt{.WithLifetime(Persistent)} & Container survives AppHost exit; reused on next \texttt{dotnet run} & Fresh empty DB on every developer session \\
\addlinespace
\texttt{.AddDatabase("gamedb")} & Registers logical DB; name must match \texttt{AddNpgsqlDbContext<>} argument & \texttt{ConnectionStrings\_\_gamedb} not injected; EF cannot resolve context \\
\addlinespace
\texttt{.WithReference(gameDb)} & Injects \texttt{ConnectionStrings\_\_gamedb} env var into web app & App starts without connection string; runtime exception on first DB call \\
\addlinespace
\texttt{.WaitFor(gameDb)} & Delays web app launch until Postgres health check passes & Race condition; app may start before DB is ready \\
\addlinespace
\texttt{.WithExternalHttpEndpoints()} & Binds HTTP/HTTPS on \texttt{localhost} outside Aspire reverse proxy & Web app only reachable through Aspire dashboard proxy \\
\bottomrule
\end{tabularx}
\end{table}

\subsection{The Postgres Resource Name and DNS Aliasing}

The string argument to \texttt{AddPostgres} --- \texttt{"postgres"} in the example --- becomes the DNS alias under which the container is reachable from other resources within the Aspire-managed network. pgAdmin uses this alias as the host when auto-configuring its server connection. The alias is also the base for the \texttt{ConnectionStrings\_\_postgres} host entry injected into dependent projects. It should be kept lowercase and DNS-safe (letters, digits, hyphens).

\subsection{The Database Resource Name and Connection String Injection}

The string argument to \texttt{AddDatabase} --- \texttt{"gamedb"} --- determines the key under which the connection string is injected into the web project's environment. Aspire injects it as \texttt{ConnectionStrings\_\_gamedb} (using double underscore as the hierarchy separator, consistent with the .NET configuration system). The \texttt{Aspire.Npgsql.EntityFrameworkCore.PostgreSQL} package reads this key automatically when \texttt{AddNpgsqlDbContext\textless AppDbContext\textgreater("gamedb")} is called in the web project's \texttt{Program.cs}. The two strings must match exactly; a mismatch produces a runtime \texttt{InvalidOperationException} with a message indicating that no connection string named \texttt{"gamedb"} was found.

\subsection{Persistent Containers and Developer Workflow}

\texttt{ContainerLifetime.Persistent} changes the container lifecycle in two specific ways. First, the container is not stopped when the AppHost process exits (\texttt{Ctrl+C} or IDE stop). Second, on the next \texttt{dotnet run}, Aspire detects the existing container by name and attaches to it rather than creating a new one. This means that any data written in session $n$ --- EF migrations, seeded rows, saved game state --- is available in session $n+1$. The developer's workflow collapses to: run, work, stop, resume; no database re-seeding is required.

The trade-off is that the container must be manually pruned when a clean state is desired. The command \texttt{docker rm -f postgres} (using the resource name) achieves this. Alternatively, removing the named volume with \texttt{docker volume rm} achieves the same effect at the data layer without removing the container.

\subsection{WaitFor and EF Migration Ordering}

\texttt{WaitFor(gameDb)} ensures the web application process does not start until the Postgres container passes its TCP health check.

\section{The ServiceDefaults Shared Library}
\label{sec:servicedefaults}

\subsection{Purpose and Project Structure}

The ServiceDefaults library is a class library that is referenced by every project the AppHost orchestrates. Its purpose is to centralise three cross-cutting concerns so that they do not need to be independently configured in each orchestrated project:

\begin{enumerate}
  \item \textbf{OpenTelemetry} --- structured logs, distributed traces, and metrics exported via OTLP to the Aspire dashboard.
  \item \textbf{HTTP resilience} --- default retry and circuit-breaker policies on outbound \texttt{HttpClient} instances via \texttt{Microsoft.Extensions.Http.Resilience}.
  \item \textbf{Service discovery} --- resolution of service names to endpoints via \texttt{Microsoft.Extensions.ServiceDiscovery}, enabling inter-service calls by name rather than hard-coded URLs.
\end{enumerate}

The \texttt{.csproj} for the HobbitGame ServiceDefaults project is:

\begin{lstlisting}[language=XML]
<!-- src/HobbitGame.ServiceDefaults/HobbitGame.ServiceDefaults.csproj -->
<Project Sdk="Microsoft.NET.Sdk">
  <PropertyGroup>
    <TargetFramework>net9.0</TargetFramework>
    <ImplicitUsings>enable</ImplicitUsings>
    <Nullable>enable</Nullable>
    <IsAspireSharedProject>true</IsAspireSharedProject>
  </PropertyGroup>

  <ItemGroup>
    <PackageReference Include="Microsoft.Extensions.Http.Resilience"    Version="9.0.0" />
    <PackageReference Include="Microsoft.Extensions.ServiceDiscovery"   Version="9.0.0" />
    <PackageReference Include="OpenTelemetry.Exporter.OpenTelemetryProtocol" Version="1.10.0" />
    <PackageReference Include="OpenTelemetry.Extensions.Hosting"        Version="1.10.0" />
    <PackageReference Include="OpenTelemetry.Instrumentation.AspNetCore" Version="1.10.1" />
    <PackageReference Include="OpenTelemetry.Instrumentation.Http"       Version="1.10.1" />
    <PackageReference Include="OpenTelemetry.Instrumentation.Runtime"    Version="1.10.0" />
  </ItemGroup>
</Project>
\end{lstlisting}

\subsection{Extensions.cs}

The library exposes two extension methods on \texttt{IHostApplicationBuilder}: \texttt{AddServiceDefaults}, called once per orchestrated project, and \texttt{MapDefaultEndpoints}, called once the application is built to register health-check routes:

\begin{lstlisting}[language={[Sharp]C}]
// src/HobbitGame.ServiceDefaults/Extensions.cs
using Microsoft.AspNetCore.Builder;
using Microsoft.AspNetCore.Diagnostics.HealthChecks;
using Microsoft.Extensions.DependencyInjection;
using Microsoft.Extensions.Diagnostics.HealthChecks;
using Microsoft.Extensions.Logging;
using OpenTelemetry;
using OpenTelemetry.Metrics;
using OpenTelemetry.Trace;

namespace Microsoft.Extensions.Hosting;

public static class Extensions
{
    public static IHostApplicationBuilder AddServiceDefaults(
        this IHostApplicationBuilder builder)
    {
        builder.ConfigureOpenTelemetry();
        builder.AddDefaultHealthChecks();
        builder.Services.AddServiceDiscovery();
        builder.Services.ConfigureHttpClientDefaults(http =>
        {
            http.AddStandardResilienceHandler();
            http.AddServiceDiscovery();
        });
        return builder;
    }

    public static IHostApplicationBuilder ConfigureOpenTelemetry(
        this IHostApplicationBuilder builder)
    {
        builder.Logging.AddOpenTelemetry(logging =>
        {
            logging.IncludeFormattedMessage = true;
            logging.IncludeScopes = true;
        });

        builder.Services.AddOpenTelemetry()
            .WithMetrics(metrics =>
            {
                metrics
                    .AddAspNetCoreInstrumentation()
                    .AddHttpClientInstrumentation()
                    .AddRuntimeInstrumentation();
            })
            .WithTracing(tracing =>
            {
                tracing
                    .AddAspNetCoreInstrumentation()
                    .AddHttpClientInstrumentation();
            });

        builder.AddOpenTelemetryExporters();
        return builder;
    }

    private static IHostApplicationBuilder AddOpenTelemetryExporters(
        this IHostApplicationBuilder builder)
    {
        var useOtlpExporter =
            !string.IsNullOrWhiteSpace(
                builder.Configuration["OTEL_EXPORTER_OTLP_ENDPOINT"]);

        if (useOtlpExporter)
        {
            builder.Services.AddOpenTelemetry().UseOtlpExporter();
        }
        return builder;
    }

    public static IHostApplicationBuilder AddDefaultHealthChecks(
        this IHostApplicationBuilder builder)
    {
        builder.Services.AddHealthChecks()
            .AddCheck("self", () => HealthCheckResult.Healthy(),
                      ["live"]);
        return builder;
    }

    public static WebApplication MapDefaultEndpoints(this WebApplication app)
    {
        app.MapHealthChecks("/health");
        app.MapHealthChecks("/alive", new HealthCheckOptions
        {
            Predicate = r => r.Tags.Contains("live")
        });
        return app;
    }
}
\end{lstlisting}

\subsection{What the Dashboard Receives}

When the AppHost is running, it starts an OTLP collector endpoint and injects \texttt{OTEL\_EXPORTER\_OTLP\_ENDPOINT} into every orchestrated project's environment. The \texttt{AddOpenTelemetryExporters} method detects this variable and activates the OTLP exporter. Logs, traces, and metrics from every orchestrated project therefore flow automatically into the Aspire dashboard without any manual configuration in the individual projects.

The \texttt{/health} and \texttt{/alive} endpoints registered by \texttt{MapDefaultEndpoints} are used by Aspire's health-check polling to determine when a resource is ready. This is the mechanism that enables \texttt{WaitFor} to function correctly for project resources (as opposed to container resources, which use TCP probing).

\section{Integrating the Web Project}
\label{sec:web}

\subsection{NuGet Package Changes}

Two additions are required in the web project's \texttt{.csproj}: a reference to \texttt{Aspire.Npgsql.EntityFrameworkCore.PostgreSQL} and a project reference to the ServiceDefaults library:

\begin{lstlisting}[language=XML]
<!-- src/Game.Web/Game.Web.csproj (additions only) -->
<ItemGroup>
  <PackageReference Include="Aspire.Npgsql.EntityFrameworkCore.PostgreSQL"
                    Version="9.0.0" />
</ItemGroup>

<ItemGroup>
  <ProjectReference Include="..\Game.ServiceDefaults\
Game.ServiceDefaults.csproj" />
</ItemGroup>
\end{lstlisting}

\texttt{Aspire.Npgsql.EntityFrameworkCore.PostgreSQL} is a thin integration package that wraps \texttt{Npgsql.EntityFrameworkCore.PostgreSQL} and adds Aspire-aware configuration reading. Specifically, its \texttt{AddNpgsqlDbContext\textless TContext\textgreater} extension method reads the connection string from the .NET configuration system under the key \texttt{ConnectionStrings\_\_\{name\}}, which is exactly the key that Aspire injects via \texttt{WithReference}. The developer does not write a connection string anywhere; Aspire provides it at runtime.

\subsection{Program.cs Changes}

Two additions are required in the web project's \texttt{Program.cs}. Both must appear before \texttt{builder.Build()}:

\begin{lstlisting}[language={[Sharp]C}]
// src/Game.Web/Program.cs (additions only; existing lines unchanged)
var builder = WebApplication.CreateBuilder(args);

// --- Aspire additions ---
builder.AddServiceDefaults();                             // (1)
builder.AddNpgsqlDbContext<AppDbContext>("gamedb");        // (2)
// --- end Aspire additions ---

// ... existing DI registrations, middleware, etc. ...

var app = builder.Build();

// --- Aspire addition ---
app.MapDefaultEndpoints();                                // (3)
// --- end Aspire addition ---

// ... existing middleware pipeline, app.Run() ...
\end{lstlisting}

Three annotations:

\textbf{(1) \texttt{AddServiceDefaults()}.} Registers OpenTelemetry, HTTP resilience, and service discovery. Must be called before any \texttt{AddHttpClient} calls so that the resilience and service-discovery handlers are applied globally.

\textbf{(2) \texttt{AddNpgsqlDbContext\textless AppDbContext\textgreater("gamedb")}.} Registers \texttt{AppDbContext} in the DI container, reads the connection string from \texttt{ConnectionStrings\_\_gamedb}, and enables Aspire health checks and telemetry for the database connection. The string \texttt{"gamedb"} must match the argument passed to \texttt{AddDatabase} in the AppHost. No \texttt{appsettings.json} entry is needed.

\textbf{(3) \texttt{MapDefaultEndpoints()}.} Registers \texttt{/health} and \texttt{/alive} endpoints. This call must appear after \texttt{builder.Build()} and before \texttt{app.Run()}.

\subsection{Avoiding Conflicts with Existing DbContext Registrations}

If the web project already registers \texttt{AppDbContext} via \texttt{AddDbContext} (for example, with a connection string read from \texttt{appsettings.json}), that registration must be removed before adding the Aspire version. Registering the same context twice produces a runtime \texttt{InvalidOperationException}. The Aspire-aware registration supersedes the manual one and does not require any compensating configuration entry in \texttt{appsettings.json}.

\subsection{Database Initialisation at Startup}

For the Game POC, schema creation uses \texttt{EnsureCreatedAsync}

\begin{lstlisting}[language={[Sharp]C}]
// Inside app.Services.CreateScope() at startup
await dbContext.Database.EnsureCreatedAsync(); // POC / development
// OR
await dbContext.Database.MigrateAsync();       // production / staged migrations
\end{lstlisting}

Both calls are safe to make under \texttt{WaitFor} ordering because Aspire guarantees the Postgres container is accepting connections before the web project process starts.

\section{Secrets and Credentials}
\label{sec:secrets}

\subsection{How Aspire Generates and Stores Passwords}

On the first \texttt{dotnet run} invocation of the AppHost, Aspire generates a random password for the PostgreSQL \texttt{postgres} superuser and stores it in the .NET User Secrets store associated with the AppHost project's \texttt{UserSecretsId}. On subsequent invocations, the stored password is read back from User Secrets and passed to the container as the \texttt{POSTGRES\_PASSWORD} environment variable. The container uses this value to initialise the database if it does not already exist, or to authenticate against an existing one.

The User Secrets store is located at:

\begin{lstlisting}
# Windows
%APPDATA%\Microsoft\UserSecrets\<UserSecretsId>\secrets.json

# macOS / Linux
~/.microsoft/usersecrets/<UserSecretsId>/secrets.json
\end{lstlisting}

The \texttt{secrets.json} file is outside the repository directory and is therefore never captured by \texttt{git add} regardless of the \texttt{.gitignore} configuration.

\subsection{Inspecting and Overriding Credentials}

To list all stored secrets for the AppHost:

\begin{lstlisting}
cd src/HobbitGame.AppHost
dotnet user-secrets list
\end{lstlisting}

A typical output for a fresh setup:

\begin{lstlisting}
Parameters:postgres-password = g7K#mP2xQn...
\end{lstlisting}

To override the generated password with a known value (useful for connecting external tools that do not read the Aspire environment injection):

\begin{lstlisting}
dotnet user-secrets set "Parameters:postgres-password" "myDevPassword"
\end{lstlisting}

The key format \texttt{Parameters:\{resource-name\}-password} is the convention Aspire uses for all generated parameter values~\cite{aspiredocs2024}.

\subsection{pgAdmin Credentials}

pgAdmin's default email and password are generated separately and appear as environment variables on the \texttt{pgadmin} resource in the Aspire dashboard. They are accessible at \texttt{PGADMIN\_DEFAULT\_EMAIL} and \texttt{PGADMIN\_DEFAULT\_PASSWORD} in the dashboard's resource-detail view. These credentials can also be overridden via \texttt{WithEnvironment} on the pgAdmin resource builder if a fixed value is preferred during development.

\subsection{Common Credential Pitfalls}

Table~\ref{tab:credential-pitfalls} documents the most common credential-related failures and their resolutions.

\begin{table}[H]
\centering
\caption{Common credential pitfalls when setting up Aspire with PostgreSQL.}
\label{tab:credential-pitfalls}
\small
\renewcommand{\arraystretch}{1.15}
\begin{tabularx}{\textwidth}{@{}>{\raggedright\arraybackslash}p{4cm} >{\raggedright\arraybackslash}X >{\raggedright\arraybackslash}p{3.5cm}@{}}
\toprule
\textbf{Symptom} & \textbf{Cause} & \textbf{Resolution} \\
\midrule
\texttt{password authentication failed for user "postgres"} & Container was created with a different password than the one now in User Secrets & \texttt{docker rm -f postgres} then re-run AppHost; new password stored \\
\addlinespace
\texttt{No connection string named "hobbitdb"} & \texttt{AddDatabase} name does not match \texttt{AddNpgsqlDbContext} name & Ensure both strings are identical \\
\addlinespace
Container starts with empty DB on every run & \texttt{WithDataVolume()} omitted or \texttt{WithLifetime(Persistent)} omitted & Add both calls to the Postgres builder chain \\
\addlinespace
pgAdmin shows no registered server & pgAdmin container was recreated after initial registration & Remove pgAdmin container; Aspire re-registers on fresh start \\
\addlinespace
Credentials appear in \texttt{appsettings.json} & Manual connection string written by developer & Remove from \texttt{appsettings.json}; rely solely on Aspire injection \\
\bottomrule
\end{tabularx}
\end{table}

\subsection{What Must Never Enter Source Control}

Three categories of data must not be committed:

\begin{itemize}
  \item \texttt{secrets.json} files from the User Secrets store (never in the repository directory, so this is enforced by location rather than \texttt{.gitignore}).
  \item Any \texttt{appsettings.\{environment\}.json} file containing a \texttt{ConnectionStrings} or \texttt{Password} key.
  \item \texttt{docker-compose.override.yml} files generated by tooling that may embed passwords.
\end{itemize}

The \texttt{.gitignore} template in this repository excludes \texttt{appsettings.*.json} by default; reviewers should verify this exclusion is present before merging a change that touches the application configuration layer.

\section{Discussion}
\label{sec:discussion}

\subsection{Scope and Limitations}

The setup described in this article is explicitly scoped to the local development environment. Aspire's AppHost orchestrator is a development-time tool: it launches processes and containers on the developer's machine, injects configuration, and provides a telemetry dashboard. It is not a production deployment mechanism. Shipping the AppHost to a server and running it there is neither the intended use case nor a supported scenario.

For production deployment, the canonical path is to use the Aspire manifest --- generated by \texttt{dotnet run --project HobbitGame.AppHost -- --publisher manifest --output-path ./manifest.json} --- as input to a deployment tool such as \texttt{azd} (Azure Developer CLI) or \texttt{aspir8} (Kubernetes manifests). These tools translate the Aspire resource graph into cloud-native deployment artifacts. The Aspire documentation covers this path in detail~\cite{aspiredocs2024}.

\subsection{EnsureCreatedAsync vs. MigrateAsync}

The HobbitGame POC uses \texttt{EnsureCreatedAsync} for schema creation. This is appropriate when the schema is stable and no migration history is maintained. It has two limitations: (i) it cannot apply incremental schema changes --- if the schema is modified after the database exists, \texttt{EnsureCreatedAsync} is a no-op and the change is not applied; (ii) it is incompatible with EF Core migrations, because it creates the schema directly from the model and does not create the \texttt{\_\_EFMigrationsHistory} table. Teams intending to evolve the schema should add a migration (\texttt{dotnet ef migrations add Init}) and switch to \texttt{MigrateAsync} at startup.

\subsection{ContainerLifetime.Persistent in Team Environments}

\texttt{ContainerLifetime.Persistent} is designed for the single-developer local workflow. In a team environment where multiple developers share a machine (uncommon) or where CI runs the AppHost (not a standard CI pattern), the persistent lifetime may produce unexpected behaviour: one developer's container state may conflict with another's, or a CI agent may accumulate stale containers across runs. For these environments, the default \texttt{ContainerLifetime.Session} --- which tears down containers on AppHost exit --- is safer.

\subsection{Generalisation to Other Databases and Services}

The pattern described for PostgreSQL generalises directly to other databases and services supported by Aspire's hosting packages. \texttt{Aspire.Hosting.Redis} provides \texttt{AddRedis}; \texttt{Aspire.Hosting.SqlServer} provides \texttt{AddSqlServer}; \texttt{Aspire.Hosting.RabbitMQ} provides \texttt{AddRabbitMQ}. All follow the same pattern: declare the resource in the AppHost, add a reference to the dependent project, call the corresponding integration package's registration method in the dependent project's \texttt{Program.cs}. The secrets, health-check ordering, and telemetry mechanisms work identically across resource types.

\subsection{Relation to Docker Compose}

Teams already using Docker Compose for local development may ask whether Aspire replaces it. The honest answer is: for the local development scenario it largely does, and with a better developer experience (typed C\# builder calls, automatic connection-string injection, integrated dashboard, no YAML). However, Compose files remain useful for (i) sharing a reproducible environment with non-.NET tooling, (ii) running the stack in CI environments where the AppHost pattern does not apply, and (iii) deploying to Docker Swarm. The two approaches can coexist: a Compose file can be maintained for CI and non-development environments while the Aspire AppHost is used for local development.

\section{Conclusion}
\label{sec:conclusion}

Adding .NET Aspire orchestration to an existing solution is, in practice, a five-step operation: create the AppHost project with its secondary SDK and two NuGet references; write a \texttt{Program.cs} that chains \texttt{AddPostgres}, \texttt{WithDataVolume}, \texttt{WithPgAdmin}, \texttt{WithLifetime(Persistent)}, \texttt{AddDatabase}, and \texttt{AddProject}; create the ServiceDefaults library and call \texttt{AddServiceDefaults()} in the web project; add \texttt{Aspire.Npgsql.EntityFrameworkCore.PostgreSQL} to the web project and register the DbContext; register both new projects in the solution file. The result is a development environment in which Postgres starts automatically, state survives restarts, pgAdmin is available at a dashboard URL, connection strings are injected without manual configuration, and OpenTelemetry signals from every project appear in a single browser dashboard.

The setup has a non-trivial number of coupling points --- the database resource name must match between AppHost and web project, the \texttt{UserSecretsId} must be set for credential persistence, \texttt{WaitFor} must precede the web project launch, and \texttt{ContainerLifetime.Persistent} must be combined with \texttt{WithDataVolume} to achieve durable state. Each of these coupling points is a potential failure mode, and Table~\ref{tab:builder-calls} and Table~\ref{tab:credential-pitfalls} are intended as quick reference for diagnosing the most common ones.

The broader observation is that the friction historically associated with local database setup --- credential management, connection-string plumbing, container lifecycle, health-check ordering --- is now automatable at the application layer in C\#, with no YAML authoring and no manual Docker commands required for the happy path. For teams whose current onboarding process includes a step labelled ``set up local Postgres,'' the Aspire AppHost pattern eliminates that step entirely.

\begin{thebibliography}{99}

\bibitem{aspiredocs2024}
Microsoft. (2024).
.NET Aspire documentation.
\textit{Microsoft Learn}.
Retrieved from \url{https://learn.microsoft.com/en-us/dotnet/aspire/}

\bibitem{microsoftusersecrets2024}
Microsoft. (2024).
Safe storage of app secrets in development in ASP.NET Core.
\textit{Microsoft Learn}.
Retrieved from \url{https://learn.microsoft.com/en-us/aspnet/core/security/app-secrets}

\bibitem{efcoredocs2024}
Microsoft. (2024).
Migrations overview --- EF Core.
\textit{Microsoft Learn}.
Retrieved from \url{https://learn.microsoft.com/en-us/ef/core/managing-schemas/migrations/}

\bibitem{otelspec2023}
OpenTelemetry Authors. (2023).
OpenTelemetry specification.
Retrieved from \url{https://opentelemetry.io/docs/specs/otel/}

\bibitem{npgsqldocs2024}
Npgsql contributors. (2024).
Npgsql Entity Framework Core provider.
Retrieved from \url{https://www.npgsql.org/efcore/}

\bibitem{pgadmindocs2024}
pgAdmin Development Team. (2024).
pgAdmin 4 documentation.
Retrieved from \url{https://www.pgadmin.org/docs/pgadmin4/latest/}

\bibitem{dockervolumes2024}
Docker Inc. (2024).
Use volumes.
\textit{Docker Documentation}.
Retrieved from \url{https://docs.docker.com/storage/volumes/}

\bibitem{aspirehosting2024}
Microsoft. (2024).
.NET Aspire PostgreSQL integration.
\textit{Microsoft Learn}.
Retrieved from \url{https://learn.microsoft.com/en-us/dotnet/aspire/database/postgresql-integration}

\bibitem{aspireservicedefaults2024}
Microsoft. (2024).
.NET Aspire service defaults.
\textit{Microsoft Learn}.
Retrieved from \url{https://learn.microsoft.com/en-us/dotnet/aspire/fundamentals/service-defaults}

\bibitem{resilience2024}
Microsoft. (2024).
Microsoft.Extensions.Http.Resilience.
\textit{NuGet Gallery}.
Retrieved from \url{https://www.nuget.org/packages/Microsoft.Extensions.Http.Resilience}

\end{thebibliography}


\end{document}
